\chapter{Теория}
\label{ch:theory}
\section{Географические координаты (LLA)}
LLA (Latitude, Longitude, Altitude) - система координат, используемая для определения положения точки на поверхности Земли. Широта - это угол между плоскостью экватора и отвесной линией в данной точке. Измеряется в градусах, изменяется от $-90^{\circ}$ (южная широта) до $+90^{\circ}$ (северная широта). Долгота - угол между плоскостью нулевого меридиана и плоскостью меридиана, проходящего через точку. Измеряется в градусах от $-180^{\circ}$ до $+180^{\circ}$. Высота - расстояние от точки до среднего уровня моря, измеряется в метрах. \cite{lla}

\section{Архитектура GSM и LTE}
\subsection{GSM (2G)}
GSM — стандарт второго поколения для голосовых вызовов, СМС и низкоскоростного интернета. Технология сохраняет актуальность благодаря множеству устройств с поддержкой только 2G \cite{gsm_lte}.

Архитектура: мобильная станция (UE), базовая станция (BTS), контроллер BSC, центр коммутации (MSC-S/MGW). Множественный доступ — TDMA/FDMA. Пакетную передачу данных обеспечивает узел Gn/Gp SGSN.


\subsection{LTE (4G)}
LTE — наиболее востребованная технология, обеспечивающая высокие скорости интернета, поддержку VoLTE и NB-IoT \cite{gsm_lte}.

Ключевые компоненты: UE, базовая станция (eNodeB), узел управления мобильностью (MME), обслуживающий шлюз (SGW), пакетный шлюз (PGW с интерфейсом Gn), сервер абонентских данных (HSS), узел политик тарификации (PCRF) \cite{gsm_lte}.


\section{Критерии качества радиосигнала}
Received Signal Strength Indicator (RSSI) — индикатор уровня принимаемого сигнала. Измеряется в дБм. Показывает общую мощность сигнала, включая полезный сигнал, шумы и помехи. Для LTE:  \cite{rssi}
\begin{itemize}
    \item Отлично: $> -65$ дБм
    \item Хорошо: $-65$ ... $-75$ дБм
    \item Удовлетворительно: $-75$ ... $-85$ дБм
    \item Плохо: $-85$ ... $-95$ дБм
    \item Нет сигнала: $\leq -95$ дБм
\end{itemize}

Reference Signal Received Power (RSRP) — мощность опорных (пилотных) сигналов от базовой станции. Измеряется в дБм. Характеризует уровень полезного сигнала без учёта шумов. Для LTE: \cite{rssi}
\begin{itemize}
    \item Отлично: $\geq -80$ дБм
    \item Хорошо: $-80$ ... $-90$ дБм
    \item Удовлетворительно: $-90$ ... $-100$ дБм
    \item Плохо: $\leq -100$ дБм
\end{itemize}

\endinput